\documentclass[11pt,letterpaper]{article}
\usepackage[T1]{fontenc}
\usepackage[utf8]{inputenc}
\usepackage[margin=0.82in,headheight=15pt]{geometry}
\usepackage{mathpazo,microtype}
\usepackage{amsmath,amssymb,booktabs,longtable,array,enumitem,xcolor,listings}
\usepackage{fancyhdr,lastpage}
\usepackage[colorlinks=true,linkcolor=teal!65!black,urlcolor=blue!55!black,
  bookmarksopen=true,pdfusetitle]{hyperref}
\hypersetup{pdftitle={An editor written in Lean},pdfauthor={EditorLean project}}
\definecolor{ink}{HTML}{183546}
\definecolor{pale}{HTML}{EFF5F7}
\setlength{\parindent}{0pt}
\setlength{\parskip}{6pt}
\setlength{\emergencystretch}{2em}
\setlist[itemize]{leftmargin=1.4em,itemsep=2pt,topsep=3pt}
\setlist[enumerate]{leftmargin=1.6em,itemsep=3pt,topsep=3pt}
\renewcommand{\arraystretch}{1.14}
\newcolumntype{P}[1]{>{\raggedright\arraybackslash}p{#1}}
\newcommand{\code}[1]{\texttt{\detokenize{#1}}}
\newcommand{\key}[1]{\textsf{#1}}
\newcommand{\note}[1]{\par\smallskip\noindent\colorbox{pale}{\parbox{\dimexpr\linewidth-2\fboxsep}{#1}}\par\smallskip}
\pagestyle{fancy}
\fancyhf{}
\fancyhead[L]{\small\textcolor{ink}{EditorLean}}
\fancyhead[R]{\small User guide, tests and verification}
\fancyfoot[L]{\small 16 November 2026}
\fancyfoot[R]{\small\thepage\ / \pageref*{LastPage}}
\renewcommand{\headrulewidth}{0.3pt}
\setcounter{tocdepth}{2}
\lstset{basicstyle=\ttfamily\small,columns=fullflexible,keepspaces=true,
  breaklines=true,frame=single,rulecolor=\color{teal!25},backgroundcolor=\color{pale},
  xleftmargin=4pt,xrightmargin=4pt,showstringspaces=false,
  literate={←}{{$\leftarrow$}}1 {→}{{$\to$}}1 {≤}{{$\le$}}1
    {≠}{{$\ne$}}1 {∧}{{$\land$}}1 {↔}{{$\leftrightarrow$}}1 {¬}{{$\neg$}}1}
\title{An editor written in Lean}
\author{EditorLean project}
\date{16 November 2026}
\begin{document}
\thispagestyle{empty}
{\small\sffamily\color{teal!65!black} PRODUCTIVITY TOOLS\par}
\vspace{12pt}
{\LARGE\bfseries\color{ink} An editor written in Lean\par}
\vspace{7pt}
{\large User guide, test catalogue and first formal proofs\par}
\vspace{5pt}
EditorLean 1.0.4 \quad | \quad Lean 4.35.0-rc3 \quad | \quad 16 November 2026
\vspace{8pt}\hrule\vspace{8pt}
\note{\textbf{Current status.} The editor works when started from Finder or with
\code{Diagnose EditorLean.command}; Finder (Launch Services) activation was
confirmed by end-to-end tests on 25--26 November 2026. Startup behaviour is
outside this document's proof claims. The requested baseline has \textbf{61 test
assertions}; the current app has three additional startup checks. This work adds
\textbf{21 universally quantified Lean theorems}, checked against the existing
editor source.}
\tableofcontents
\vfill
{\small The contents entries and page numbers above are clickable. This PDF is
compiled from \code{EditorLean.tex}; the proof source and reproducible build commands
are described in Section~\ref{sec:reproduce}.}
\clearpage

\section{User guide}\label{sec:guide}
\subsection{Starting and using the editor}
EditorLean is a small NEdit-style editor for UTF-8 text files. Its editor logic is
written in Lean, compiled to C with \code{lean -c}, and then compiled to a native
Apple-silicon executable with \code{leanc}. Tcl/Tk supplies the desktop widgets
through a small C bridge. The packaged app does not need Python, Lean or Homebrew
to be installed in order to run. This build targets macOS 14 or later.

\begin{enumerate}
\item In Finder, open the project's \code{dist} folder and double-click
\code{Diagnose EditorLean.command}. Keep the resulting Terminal window open while
using the editor. It records startup diagnostics in \code{EditorLean-startup.log}.
\item Choose \key{File > Open} to open a text file, or type in the empty document.
A leading asterisk in the window title means the text differs from its saved state.
\item Use \key{File > Save} or \key{Cmd+S}. A new document asks for a filename.
\key{Save As} chooses a new destination.
\item Use \key{Search > Find / Replace} for literal, case-sensitive searches.
The search wraps at the beginning or end of the document.
\item Quit normally when finished, saving or discarding changes when asked.
After an app update, quit the old running editor before launching the new build.
\end{enumerate}

Each document has its own window. The status line shows the cursor's line and
column, insert/overtype/recording mode, UTF-8 encoding and file path. Split panes
are two views of the \emph{same} document; edits and undo history are shared.

\subsection{Menu options}
\begin{longtable}{@{}P{0.15\linewidth}P{0.79\linewidth}@{}}
\toprule\textbf{Menu} & \textbf{What it does}\\\midrule\endfirsthead
\toprule\textbf{Menu} & \textbf{What it does}\ (continued)\\\midrule\endhead
File & New; Open; Open Selected (use selected text as a path); Include File (insert
another file's text); Save; Save As; Revert; Close; Print; Quit. Revert reloads the
disk version after confirmation if there are unsaved edits. Print opens a local
HTML preview in your browser; print there with Cmd+P.\\
Edit & Undo/Redo; Cut/Copy/Paste; Paste Column; Select All; Shift Left/Right;
Uppercase/Lowercase; Fill Paragraph. Column paste inserts successive clipboard
lines at the same column, padding short lines. Case conversion uses the selection,
or the character at the cursor if there is no selection.\\
Search & Find/Replace; Find Next/Previous; Find Selection; Incremental Search;
Go to Line; Set/Go to Mark; Find Matching; Find Definition; Call Tip. Bookmarks
use names a--z. Matching handles nested bracket pairs and simple quote delimiters;
it is not a programming-language parser. Definition and Call Tip read a ctags file.\\
Preferences & Toggle line wrapping or the status line; increase/decrease the font
size. Wrapping changes the display, not the saved line breaks.\\
Shell & Execute Command; Filter Selection; Execute Current Line; Cancel.
Commands run through \code{/bin/sh}. Filter Selection supplies the selected text
as input. A filter replaces its original selection only if the document still
matches its pre-command text; otherwise output opens in a new window.\\
Macro & Learn/Finish records editor actions; Replay repeats the recorded actions;
Repeat repeats the last action a requested number of times. Cancel stops a replay
or a running shell command.\\
Windows & Split Pane; Close Pane; Next Document; Previous Document. Closing a
pane keeps the document. Closing a document can prompt to save.\\
Help & Keyboard Shortcuts lists the installed bindings. About reports the app
and toolchain versions.\\
\bottomrule
\end{longtable}

\subsection{Keyboard shortcuts}
\textbf{Notation.} The application installs both Control and Command aliases for
its Control bindings. The tables use \key{Cmd} for convenience; use \key{Ctrl}
if macOS reserves a Command combination. \key{Option} is the key called
\key{Alt} in the binding list. Some keys, such as Insert and Page Up, require an
extended keyboard or its equivalent. The Help window lists the complete bindings,
including control-character insertion and selection variants.

\begin{longtable}{@{}P{0.34\linewidth}P{0.60\linewidth}@{}}
\toprule\textbf{Keys} & \textbf{Action}\\\midrule\endfirsthead
\toprule\textbf{Keys} & \textbf{Action}\ (continued)\\\midrule\endhead
Cmd+N / O / S & New / Open / Save.\\
Shift+Cmd+S & Save As.\\
Cmd+W / Q & Close document / Quit application.\\
Cmd+Y / Option+I & Open Selected / Include File.\\
Shift+Cmd+E / Cmd+P & Revert / Print via browser preview.\\
Cmd+Z / Shift+Cmd+Z & Undo / Redo.\\
Cmd+X / C / V / A & Cut / Copy / Paste / Select All.\\
Shift+Cmd+V & Paste Column.\\
Cmd+9 / Cmd+0 & Shift selected lines left / right by four spaces.\\
Cmd+6 / Shift+Cmd+6 & Uppercase / Lowercase (the shifted binding uses the caret keysym).\\
Cmd+J & Fill Paragraph to a target width of 80 characters. Long words are not split.\\
Return / Cmd+Return & Insert newline / Insert newline with the current leading indentation.\\
Insert & Toggle overtype mode.\\
Cmd+Backspace & Delete back to a word boundary.\\
Cmd+U / Cmd+Delete & Delete to line start / line end.\\
\midrule
Cmd+F / Cmd+R & Find / Replace dialog.\\
Cmd+G / Shift+Cmd+G & Next / Previous match.\\
Cmd+H / Shift+Cmd+H & Find Selection forward / backward.\\
Cmd+I / Shift+Cmd+I & Incremental search forward / backward.\\
Cmd+T / Shift+Cmd+T & Replace and find the next / previous occurrence.\\
Option+T / Shift+Option+T & Replace forward / backward without the additional find step.\\
Cmd+L / Cmd+E & Go to Line / Go to the selected line or line:column.\\
Option+M / Option+G & Set / Go to a named mark. Add Shift to Go to Mark to extend selection.\\
Cmd+M / Shift+Cmd+M & Find matching delimiter / Extend selection to it.\\
Cmd+D / Cmd+apostrophe & Find Definition / Call Tip using ctags.\\
\midrule
Arrows; Home / End & Move; line start / end. Shift extends the selection.\\
Ctrl+Left / Right & Move by words.\\
Ctrl+Up / Down & Move toward a paragraph boundary.\\
Ctrl+Home / End & Document start / end.\\
Option+Shift+arrows & Extend a rectangular selection.\\
Ctrl+Space & Set a selection anchor.\\
Ctrl+Shift+Space & Extend from the anchor; add Option for a rectangle.\\
Cmd+2 / Cmd+1 & Split / Close Pane.\\
Cmd+Page Down / Up & Next / Previous document.\\
Option+Home & Switch to the previously active document.\\
Cmd+plus / minus & Increase / Decrease font size.\\
Option+A & Show / Hide status line.\\
Option+X / Option+R & Execute shell command / Filter Selection.\\
Ctrl+keypad Enter & Execute Current Line.\\
Option+K / Cmd+K & Learn/Finish macro / Replay.\\
Cmd+comma & Repeat Last Action.\\
Escape / Cmd+period & Cancel current operation.\\
F1 / F10 & Shortcut help / Open the File menu.\\
\bottomrule
\end{longtable}

\subsection{Files, encodings and practical limits}
Only UTF-8 text is accepted. An initial UTF-8 byte-order mark (BOM) is remembered.
Line endings are normalized to LF in memory. If the text is edited, saving uses
the original file's \emph{first detected} newline style. If the text is unchanged,
the original bytes are reused exactly, including mixed line endings and the
presence or absence of a final newline.

The editor polls named files for external changes and offers to reload them.
Saving checks the destination bytes against the expected saved bytes, asks before
overwriting a conflict, stages the new data beside the destination, then renames
it over the destination. This is not a lock against another writer: a file can
change after the comparison and before the rename. Symlink targets and ordinary
permission bits are covered by integration tests, not by the current proofs.

Find is literal and case-sensitive, not a regular-expression search. Ctags
addresses support numeric lines and a limited literal-line form. Bookmarks and
macros are session state. The application is a text editor, not a syntax-aware IDE.

\clearpage
\section{Unit tests}\label{sec:tests}
\subsection{Scope and execution}
The requested \textbf{61 tests} are the 61 successful assertion instances in the
1.0.3 baseline, generated by one \code{selfTests} routine and its loops. They are
not 61 separate test functions. IDs below refer to that baseline. Version 1.0.4
adds three assertions for noninteractive Finder startup, so the current command
reports \textbf{64 passed}. Those additions are not silently included in the 61.

\begin{lstlisting}
dist/EditorLean.app/Contents/MacOS/EditorLean --self-test
\end{lstlisting}

The suite runs compiled Lean and the real Tcl/C bridge, uses temporary files for
file tests, and launches small controlled shell commands. It does not open GUI
windows. A separate \code{--gui-self-test} suite exercises real widgets,
key-event bindings, search, undo, clipboard, split panes, rectangles, macros and
shell filtering. All 20 GUI tests passed in a normal desktop session on
26 November 2026, including runs without a terminal attached.

\subsection{The 61-test catalogue}
The seven groups contain $2+4+21+9+9+13+3=61$ assertions. The interop group includes
ID~27; this explains the order around the keyboard cases.

{\fontsize{9.5}{11}\selectfont\renewcommand{\arraystretch}{1.04}
\begin{longtable}{@{}P{0.09\linewidth}P{0.28\linewidth}P{0.57\linewidth}@{}}
\toprule\textbf{ID} & \textbf{Case} & \textbf{What is checked}\\\midrule\endfirsthead
\toprule\textbf{ID} & \textbf{Case} & \textbf{What is checked}\ (continued)\\\midrule\endhead
\multicolumn{3}{@{}l}{\textbf{Startup: 2 assertions}}\\*
1 & Main thread & Native application entry remains on the macOS main thread.\\
2 & Worker guard & A worker-thread attempt to initialize the GUI is rejected before AppKit is entered.\\
\midrule
\multicolumn{3}{@{}l}{\textbf{Tcl/C interoperation: 4 assertions}}\\*
3 & String round trip & A sample containing Greek, emoji, CJK, punctuation, newlines and embedded NUL survives a Tcl variable round trip.\\
4 & List quoting & An array containing empty, complex, spaced and newline strings survives Tcl list construction and parsing.\\
5 & Callback arguments & A Tcl command invokes exported Lean code with the original argument array intact.\\
27 & Error propagation & An unknown Tcl command becomes an IO error across the C bridge.\\
\midrule
\multicolumn{3}{@{}l}{\textbf{Keyboard handling: 21 assertions}}\\*
6--21 & Literal insertion & One assertion for each input: l, a, 1, space, double quote, apostrophe, left/right brace, left/right bracket, dollar sign, semicolon, backslash, Greek alpha, emoji, and the text \code{[set injected 1]}. Each is delivered literally to the Lean callback; the handler returns Tcl's break code to prevent duplicate default insertion.\\
22--25 & Ignored events & Four cases produce no insertion and a normal return: empty text with state 0, and x with modifier states 4, 8 and 12.\\
26 & No evaluation & The command-like text from ID~21 did not create the Tcl variable \code{injected}.\\
\midrule
\multicolumn{3}{@{}l}{\textbf{Pure editing helpers: 9 assertions}}\\*
28 & Nested forward match & The opening parenthesis in \code{a(b(c)d)e} matches offset 7.\\
29 & Nested backward match & The closing parenthesis at offset 7 matches offset 1.\\
30 & Unmatched delimiter & An unmatched opening parenthesis returns no match.\\
31 & Quotes & The first quote in a quoted sample matches the closing quote.\\
32 & Word boundaries & Forward and backward navigation on an ASCII two-word sample gives the expected offsets.\\
33 & Unicode boundaries & Greek and CJK text yields the expected word offsets. Classification calls Tcl's Unicode character predicate.\\
34 & Indentation & Adding two spaces indents both lines, retaining existing indentation.\\
35 & Unindentation & Removing four spaces handles an indented line, a leading tab and a line with fewer than four spaces.\\
36 & Paragraph fill & Extra spaces and a line break collapse into one indented paragraph.\\
\midrule
\multicolumn{3}{@{}l}{\textbf{Encoding: 9 assertions}}\\*
37 & Unchanged plain text & Plain text without a final newline retains its exact bytes.\\
38 & Unchanged LF & A trailing LF is preserved exactly.\\
39 & Unchanged CRLF & CRLF bytes survive an unchanged decode/encode round trip.\\
40 & Unchanged CR & A CR-only newline survives the round trip.\\
41 & Unchanged BOM/Unicode & BOM, Greek letters and CRLF bytes are all preserved.\\
42 & Unchanged mixed endings & A file mixing CRLF, LF and CR retains every original byte.\\
43 & Unchanged empty file & The empty byte sequence stays empty.\\
44 & Edited BOM/CRLF & Replacing the text preserves the BOM and emits CRLF newlines.\\
45 & Invalid UTF-8 & Bytes FF FE are rejected instead of being decoded as text.\\
\midrule
\multicolumn{3}{@{}l}{\textbf{File integration: 13 assertions}}\\*
46 & Open normalization & A temporary CRLF file opens with LF in the buffer.\\
47 & Save format & Saving edited text writes the remembered CRLF format. This checks resulting bytes, not crash-atomicity under every failure.\\
48 & External conflict & A save after an external edit raises a conflict.\\
49 & Preserve external edit & The rejected save leaves the externally written bytes in place.\\
50 & Explicit overwrite & An explicitly authorized overwrite writes the requested text.\\
51 & Deleted destination & Deleting a previously saved destination causes an unforced save to fail.\\
52 & New document & Opening a nonexistent path produces empty text with no original bytes.\\
53 & Save As collision & Saving an unrelated new document over an existing file is blocked without overwrite authorization.\\
54 & Invalid destination & A write into an invalid path fails.\\
55 & Failed-write preservation & The earlier file still contains its original data after that failed write.\\
56 & Symlink target & Saving through a symlink updates the resolved target.\\
57 & Symlink retained & The symlink remains a symlink after the save.\\
58 & Permissions & Mode 0640 survives a save.\\
\midrule
\multicolumn{3}{@{}l}{\textbf{Asynchronous shell integration: 3 assertions}}\\*
59 & Concurrent pipe I/O & A command echoes 200,000 input characters, emits stderr and exits successfully without a pipe deadlock.\\
60 & Cancellation & Cancelling a controlled sleep job completes the task after process-group termination. This is not an exhaustive test of every possible descendant process.\\
61 & Failure status & A controlled failing command returns exit code 7 and the expected stderr.\\
\bottomrule
\end{longtable}}

\note{\textbf{What passing means.} These examples exercise the real executable and
catch regressions. They do not establish a result for every file, every Unicode
string, every schedule or every possible filesystem failure. The three additional
startup assertions check the native startup marker, noninteractive mode and empty
Tcl command-line arguments; they do not prove that Finder launches successfully.}

\clearpage
\section{Formal verification}\label{sec:formal}
\subsection{From examples to universal specifications}
A useful specification quantifies over \emph{all inputs in a stated domain} and
makes its preconditions explicit. For example, an unchanged document with original
bytes should encode to those exact bytes, regardless of their size or line endings.
That is stronger than seven selected examples. Conversely, ``every shell command
finishes'' is false: some commands deliberately never terminate.

There are three distinct scopes in this document:
\begin{itemize}
\item \textbf{Actual implementation:} theorems refer to functions imported directly
from \code{EditorLean.lean}, such as \code{encodeDocument} and \code{charSlice}.
\item \textbf{Connected save-policy abstraction:} \code{checkedSave} names the
existing post-read save body. A checked definitional equality connects it to the
actual \code{saveDocument}; it is not an unrelated idealized algorithm.
\item \textbf{Proposed future specifications:} GUI, Tcl semantics, system calls and
process scheduling require models or explicit interface assumptions. No such
assumptions have been disguised as completed application proofs here.
\end{itemize}

\subsection{Priority ranking}\label{sec:ranking}
Priority balances the value of preventing data loss against expected proof effort.
Ease estimates concern proving meaningful behavior of this implementation, not
merely writing an attractive specification. Baseline IDs are from
Section~\ref{sec:tests}; related properties may overlap.

{\small
\begin{longtable}{@{}P{0.05\linewidth}P{0.28\linewidth}P{0.16\linewidth}P{0.43\linewidth}@{}}
\toprule\textbf{Rank} & \textbf{Property / baseline IDs} & \textbf{Value; ease} & \textbf{Status and scope}\\\midrule\endfirsthead
\toprule\textbf{Rank} & \textbf{Property / baseline IDs} & \textbf{Value; ease} & \textbf{Status and scope}\ (continued)\\\midrule\endhead
1 & Unchanged bytes, 37--43 & Very high; easy & \textbf{Proved.} Any original byte array is returned unchanged when encoding the saved text. This is an encoder theorem, not yet a full decoder/encoder round-trip theorem.\\
2 & Save conflicts, 48--51, 53 & Very high; easy--moderate & \textbf{Proved decision and rejection branch.} Any mismatch is rejected unless overwrite is explicit; existing Save As targets and deleted files are covered. Disk behavior and concurrent changes remain outside the proof.\\
3 & BOM and edited output, 44 & High; easy--moderate & \textbf{Proved.} Edited BOM documents have prefix EF BB BF and the specified encoded body. The complete semantics of newline replacement are not proved.\\
4 & Decoder checks / new document, 45, 46, 52 & High; moderate & \textbf{Partly proved.} The UTF-8 primitive rejects every invalid byte array; valid strings round-trip through UTF-8; the absent-file decoder case is empty. BOM removal and newline normalization in the full decoder remain open.\\
5 & Indent/unindent, 34--35 & Medium; moderate & Candidate: positive indentation followed by matching unindentation restores all lines. Must state tab behavior and exclude or define zero-width behavior. Not proved.\\
6 & Navigation and matching, 28--33 & High; moderate--hard & Candidates: returned offsets stay in bounds; word loops terminate and move in the requested direction; bracket matches obey nesting. Requires loop invariants and a Unicode-classification contract. Only auxiliary slicing bounds are proved.\\
7 & Literal typing, 6--26 & High; hard & Candidate: arbitrary text is passed as data, never evaluated as Tcl, and inserted once. Needs a model of Tk event substitution, the Tcl procedure and the callback boundary. Current protection is tested, not formally proved.\\
8 & Paragraph fill, 36 & Medium; moderate--hard & Candidate: preserve word order and paragraph structure. A blanket 80-character bound is false for long words or indentation; a useful width theorem needs explicit preconditions. Not proved.\\
9 & File operations, 47, 54--58; disk part of 49--51 & Very high; hard & Candidates: failed staging never truncates the destination; rename behavior, symlinks and permissions satisfy contracts. Requires filesystem/system-call semantics and concurrency assumptions. Integration tests only.\\
10 & Tcl transport, 3--5, 27 & High; hard & Candidates: all well-formed strings and lists survive the ABI with their lengths, NULs and errors intact. Requires C/Lean representation and Tcl contracts. Integration tests only.\\
11 & Threading, 1--2 & Medium; hard & Candidate: all GUI operations execute on the OS main thread, not just initialization. Requires a model of the native launcher and event dispatch. Runtime guards and tests only.\\
12 & Shell jobs, 59--61 & Medium; very hard & Candidates: faithfully report finite output/exit status; eventual completion under termination and scheduling assumptions; cancellation reaches the intended group. These are temporal/OS properties, not pure Lean function equalities.\\
\bottomrule
\end{longtable}}

\subsection{Proofs obtained}\label{sec:proved}
The proof file \code{editor_lean/EditorLeanProofs.lean} imports the real application
module. It contains \textbf{21 public universal theorems}: two preservation
results, eight save-policy/body results, four edited-encoding results, three
decoder-primitive/empty-document results and four slicing results. One additional
private lemma establishes lawful byte-array equality for the save proofs.

\subsubsection*{1. Exact byte preservation}
\begin{lstlisting}
theorem unchanged_bytes (doc : Document) (raw : ByteArray)
    (h : doc.original = some raw) :
    encodeDocument doc doc.text = raw := by
  simp [encodeDocument, h]
\end{lstlisting}
This quantifies over \emph{every} document and byte array satisfying the stated
original-bytes condition. No UTF-8 validity assumption is needed for this encoder
branch: it reuses the original bytes. The second theorem,
\code{unchanged_without_original}, proves that the unchanged branch uses the text's
UTF-8 encoding when no original bytes exist. Neither theorem independently proves
that the decoder populated the document record correctly.

\subsubsection*{2. Save-policy correctness and a write-free rejection branch}
Let $E$ be the expected bytes: the saved original bytes when the canonical target
path equals the document path, and ``absent'' otherwise. The checked equivalence is
\[
  \text{conflict} = \text{true}
  \quad\Longleftrightarrow\quad
  (\text{overwrite} = \text{false})\ \land\ (\text{existing} \ne E).
\]
For arbitrary byte arrays, the proofs cover a changed same-path file, an existing
Save As destination, a deleted previously saved file, a new absent file, and an
explicit overwrite. The comparison is against the bytes \emph{already read}; it
cannot prevent a later race with another writer.

The connection to the real implementation is itself checked:
\begin{lstlisting}
theorem saveDocument_decomposition
    (doc : Document) (text path : String) (overwrite : Bool) :
    saveDocument doc text path overwrite = (do
      let path <- canonical path
      let existing <- readOptional path
      checkedSave doc text path existing overwrite) := by
  rfl
\end{lstlisting}
Here the listing writes Lean's left arrow as ASCII \code{<-}. In the source it is
spelled with the equivalent Lean arrow notation. The proof is definitional: the
named body unfolds to the existing application code without changing it.

The theorem \code{rejected_save_is_error} establishes, for all its inputs, that
when the conflict flag is true, \code{checkedSave} is equal as an IO computation
to throwing the conflict error. The later \code{atomicWrite} operation is absent
from that branch. This is stronger than proving the Boolean flag alone, but it is
\emph{not} a proof of the native write routine, nor of atomic compare-and-swap on
a live filesystem.

\subsubsection*{3. Edited encoding and the BOM prefix}
For any changed text and a document with its BOM flag set, the result is
\[
  [239,187,191]\;{+\!+}\;
  \operatorname{UTF8}\bigl(\operatorname{replaceLF}(\text{text},\text{newline})\bigr).
\]
The numbers are the actual UTF-8 bytes EF BB BF, not just a textual marker.
\begin{lstlisting}
theorem changed_bom_prefix (doc : Document) (text : String)
    (htext : text ≠ doc.text) (hbom : doc.bom = true) :
    (encodeDocument doc text).extract 0 3 =
      ByteArray.mk #[239,187,191] := by
  rw [changed_bom_encoding doc text htext hbom]
  exact ByteArray.extract_append_eq_left (by rfl)
\end{lstlisting}
Other theorems establish the full encoded-body formula, the formula with no
inserted BOM when the flag is false, and a size of three bytes plus the encoded
body length. ``No inserted BOM'' does not mean the user's text cannot itself
begin with a BOM character. The changed-text hypothesis matters: unchanged text
takes the original-byte-preserving branch instead.

\subsubsection*{4. Decoder primitives and absent documents}
\begin{lstlisting}
theorem invalid_utf8_rejected (bytes : ByteArray)
    (h : ¬ bytes.IsValidUTF8) :
    String.fromUTF8? bytes = none := by
  simp [String.fromUTF8?, h]
\end{lstlisting}
The source also proves that encoding any Lean string and decoding it with
\code{String.fromUTF8?} returns that string. These are properties of the exact
primitive used by the editor; they do not yet lift through the full BOM/newline
processing pipeline. For every path, \code{absent_document_is_empty} proves that
\code{decodeDocument path none} returns the empty document record with that path.
It does not prove the operating system accurately reports absence.

\subsubsection*{5. Bounds for the slicing helper}
For all strings $s$ and natural-number endpoints $a,b$, the proof establishes
\[
 \operatorname{length}(\operatorname{charSlice}(s,a,b))
 = \min\bigl(b\mathbin{-}a,\;\operatorname{length}(s)\mathbin{-}a\bigr),
\]
where natural-number subtraction saturates at zero. Corollaries show the slice
cannot exceed the requested length or the original length, and is empty when
$b\le a$. The bounds count Unicode scalar characters, not UTF-8 bytes or visual
grapheme clusters. They concern the actual helper; they are not a proof of the
separate word-boundary loop or every cursor operation.

\subsection{Proof search and trust boundaries}\label{sec:trust}
The proof search used the installed Lean 4.35.0-rc3 standard library and official
Lean documentation. Useful ingredients were simplification with \code{simp},
definitional equality with \code{rfl}, arithmetic reasoning with \code{omega},
string/byte-array append lemmas and list take/drop length lemmas. No Mathlib
dependency was needed. The scripts are new proofs of this editor's properties,
not purported pre-existing proofs of the whole application.

The audit command \texttt{\#print axioms} was run for every public theorem. Its
reported dependencies are subsets of \code{propext}, \code{Classical.choice} and
\code{Quot.sound}, Lean's standard logical axioms. No result uses \code{sorry},
\code{admit}, a new application axiom, \code{native_decide}, or
\code{Lean.trustCompiler}. The reproducible checker rejects theorem audits
containing axioms outside that list.

\textbf{This is not a proof of the entire executable.} The kernel checks statements
about Lean definitions. Runtime correspondence still relies on Lean's compiler
and runtime, compiled implementations of standard primitives, the handwritten C
bridge, Tcl/Tk, and macOS. The save-body equality establishes program structure;
it does not supply a formal filesystem. The certificate does not promise that
Finder works, that every GUI action is correct, or that arbitrary concurrent
writes cannot race.

\subsection{Open properties and counterexamples}
Two proposed-looking specifications need correction before attempting proofs:
\begin{itemize}
\item \textbf{``Indenting by zero is always the identity'' is false.} Evaluating
\code{indentText} on a tab followed by x, with amount zero, produces x. The
nonpositive branch removes a leading tab. A future round-trip theorem should
specify positive indentation widths and the intended tab policy.
\item \textbf{``Every filled line has length at most 80'' is false.} Evaluating
\code{fillText} on an 81-character word leaves length 81. Long words are not split.
A width theorem needs assumptions on word length and indentation, or a deliberate
behavior change.
\end{itemize}
These evaluations are recorded in \code{VerificationExamples.lean}. They are
\emph{examples}, not universal proofs or assumptions used by the theorem file.
The attempted direct \code{decide} reduction also encountered opaque string
operations; native evaluation was used only to inspect examples, never to certify
a theorem.

The next valuable work is to prove the complete decoder/encoder connection,
followed by positive-width indentation round trips and navigation invariants.
The existing partial/loop-heavy implementations and older string operations may
need proof-friendly lemmas or carefully justified refactoring. The Tcl keyboard
handler requires a semantic model before its example tests can become an
end-to-end universal non-evaluation theorem.

\subsection{Reproducing and maintaining this document}\label{sec:reproduce}
From the project directory:
\begin{lstlisting}
python3 editor_lean/check_proofs.py
python3 editor_lean/build_document.py
\end{lstlisting}
The proof checker first recompiles the current \code{EditorLean.lean} to an importable
module, then checks \code{EditorLeanProofs.lean} and audits its axioms. It writes
\code{verification.log} and \code{verification.json}, including source SHA-256
hashes. It does not modify the application or rebuild its app bundle.

The document builder rechecks the proofs and runs LaTeX until the table of contents
and cross-references are current. The final files are \code{EditorLean.tex} and
\code{EditorLean.pdf} in the project root. Re-run the builder after every document
revision; subsequent work on this document should deliver a newly compiled PDF
alongside the updated source.

\textbf{Evidence for this edition.} The current native suite reported 64 passes;
the 61 baseline assertions are catalogued above. The checker accepted 21 public
universal theorems. The application source digest is:
{\small\begin{quote}\path{2063a8b026d2b3a19c4689cbcc091d528962d9c7232d839cbfa377862aef5b39}\end{quote}}
The full audit, proof digest and exact toolchain version are in
\code{editor_lean/verification.json}. These records bind this proof run to its
source; a later source change requires another check.

\subsubsection*{Primary references}
\begin{itemize}
\item Local implementation: \path{editor_lean/EditorLean.lean}.\\
Proof module: \path{editor_lean/EditorLeanProofs.lean}.\\
Audit: \path{editor_lean/verification.log}.
\item \href{https://lean-lang.org/doc/reference/latest/Tactic-Proofs/Tactic-Reference/}{Lean Reference: Tactic Reference},
and \href{https://lean-lang.org/doc/reference/latest/ValidatingProofs/}{Validating Proofs}.
\item \href{https://lean-lang.org/doc/reference/latest/Axioms/}{Lean Reference: Axioms}.
These explain kernel checking and why native evaluation changes the trust assumptions.
\item Installed version's \href{https://github.com/leanprover/lean4/blob/v4.35.0-rc3/src/lean/Init/Data/ByteArray/Lemmas.lean}{ByteArray lemmas}
and \href{https://github.com/leanprover/lean4/blob/v4.35.0-rc3/src/lean/Init/Data/String/Defs.lean}{String definitions}.
The pinned local compiler, rather than the changing ``latest'' manual, checked the proofs.
\end{itemize}
\end{document}
